\documentclass[10pt,a4paper]{amsart}
\usepackage[margin=19mm]{geometry}
\usepackage{amsmath,amssymb,mathtools}
\usepackage[scr=rsfs,cal=euler]{mathalfa}
\usepackage{xfrac}
\usepackage[unicode,hidelinks,bookmarks=false]{hyperref}
\newcommand{\ie}{i.e.\ }
\newcommand{\cf}{cf.\ }
\newcommand{\resp}{respectively\ }
\newcommand{\Subsection}{\subsection}
\providecommand{\nobreakdash}{\nobreak\hbox{-}}
\setlength{\emergencystretch}{2em}
\multlinegap0pt
\usepackage{bm}
\usepackage{bbm}
\usepackage{dsfont}
\usepackage{xspace}
\usepackage{cancel}
\usepackage{mathtools}
\usepackage{amsthm}
\makeatletter
\@ifundefined{theorem}{\newtheorem{theorem}{Theorem}}{}
\@ifundefined{lemma}{\newtheorem{lemma}[theorem]{Lemma}}{}
\makeatother
\newcommand*\diff{\mathop{}\!\mathrm{d}}
\title[Neural Hamilton--Jacobi Characteristic Flows for Optimal Tra]{Neural Hamilton--Jacobi Characteristic Flows for Optimal Transport}
\author{Yesom Park, Shu Liu, Mo Zhou, Stanley Osher}
\date{Source: arXiv:2510.01153v1 (CC BY 4.0)}
\begin{document}
\maketitle

\noindent\emph{Source and licence.} Neural Hamilton--Jacobi Characteristic Flows for Optimal Transport. Version arXiv:2510.01153v1 (CC BY 4.0), distributed under the Creative Commons Attribution 4.0 International licence, as recorded in the arXiv licence metadata for this exact version and shown on the abstract page https://arxiv.org/abs/2510.01153v1. Full rights notice: licenses/arxiv-2510.01153-CC-BY-4.0.txt.\par\smallskip

\noindent\textit{Editorial scope.} Counted unit: the Theorem whose source label is \texttt{thm:main} in the article (source0.tex lines 712--719, source label \texttt{thm:main}), delivered together with the article's own complete original proof of it (source0.tex lines 720--757, one \texttt{proof} environment of 38 lines). No mathematics is written, repaired or re-derived: statement and proof are the article's own text line for line. Required context retained uncounted: lemma:sufficient\_cond\_for\_OT (lines 698--700); lemma:mono\_map (lines 703--705). Every definition, display and label of those lines is retained unchanged. Omitted from this package and not redistributed: 1--697, 701--702, 706--711, 758--1708. Nothing inside a retained excerpt is omitted or altered: every retained body is delivered exactly as the article's lines are written, so no omission receipt of any kind accompanies this package. Citations: the retained excerpts carry no citation command at all, so no bibliography entry of the article is transcribed and no reference is redistributed. Reference adaptation: none was needed and none was made; every reference inside the delivered unit resolves inside it, so no unresolved reference or citation is produced. Printed slips kept literal: no printed word, symbol, hypothesis or number of the counted statement or of its proof was altered. Layout and class: the article's own class is \texttt{article} with packages including iclr2025\_arXiv, times, bbm, hyperref, url, tikz, multirow, enumitem, subfigure, booktabs, subcaption, algorithm, algorithmicx, algpseudocode; the wrapper is the lane's pinned \texttt{amsart} editorial wrapper at 10pt on A4, which restates the theorem-like environments the retained text needs and adds the article's own macro definitions that the retained text needs (\texttt{\textbackslash diff}), with extra wrapper packages (bm, bbm, dsfont, xspace, cancel, mathtools). The delivered numbering is the unit's own. Assets: no retained span contains a figure environment, a graphics-inclusion command, a PDF-page inclusion, a table or a TikZ picture; no image file is copied into this package, the package declares no asset dependency and no image file is redistributed with it. Licence: version arXiv:2510.01153v1 of this article is distributed under the Creative Commons Attribution 4.0 International licence, as recorded in the arXiv licence metadata for this exact version and shown on the abstract page https://arxiv.org/abs/2510.01153v1. A full rights notice is delivered at licenses/arxiv-2510.01153-CC-BY-4.0.txt. Characters: every retained excerpt is pure ASCII, so the delivered engine, XeLaTeX with its default Latin Modern text font, reports no missing character.
\begin{lemma}\label{lemma:sufficient_cond_for_OT} 
  Suppose that $\mu, \nu$ are probability distributions on $\mathbb{R}^d$. Suppose $\mu$ has finite second order moment, i.e.,  $\int_{\mathbb{R}^d} \|x\|^2 d\mu(x) < \infty$. Assume that $\psi:\mathbb{R}^d \rightarrow \mathbb{R}\bigcup \{+\infty \}$ is convex and differentiable $\mu$-a.e. Set $T = \nabla \psi$ and suppose $\int_{\mathbb{R}^d} \|T(x)\|^2 d\mu(x) < + \infty$. Then $T$ is optimal for the transport cost $\frac12 \|x - y\|^2$ between $\mu$ and $\nu$. 
\end{lemma}

\begin{lemma}\label{lemma:mono_map}
  Suppose $T:\mathbb{R}^d\rightarrow\mathbb{R}^d$ is a bijective map on $\mathbb{R}^d$. Assume that $T$ is strictly monotone, that is, $\langle T(x) - T(y), x-y \rangle > 0$ for arbitrary $x, y\in\mathbb{R}^d$, $x\neq y$. Here, we denote $\langle\cdot, \cdot \rangle$ as the $\ell^2$ inner product on $\mathbb{R}^d$. Suppose $S:\mathbb{R}^d\rightarrow \mathbb{R}^d$ satisfies $S\circ T = \mathrm{Id}$, then $S$ is also bijective, and strictly monotone.
\end{lemma}

\begin{theorem}\label{thm:main}
  Given the probability distributions $\mu, \nu \in \mathscr{P}(\mathbb{R}^d)$ with $\int_{\mathbb{R}^d}\|x\|^2\diff \mu, \int_{\mathbb{R}^d}\|x\|^2\diff \nu < +\infty,$ 
  suppose $u_0 \in C^1_{\text{loc}}(\mathbb{R}^d), u_1 \in C^2_{\text{loc}}(\mathbb{R}^d), \nabla u_1 \in L^2(\mathbb{R}^d, \mathbb{R}^d; \nu)$ satisfy
  \begin{equation}
    u_0(x-t_f\nabla u_1(x)) = u_1(x) - \frac{t_f}{2}\|\nabla u_1(x)\|^2, \quad \forall \ x\in\mathbb{R}^d.  \label{implicit_equ_at_T}
  \end{equation}
  Assume further that $(\mathrm{Id} + t_f\nabla u_0(\cdot))_\sharp\mu = \nu$. If the mapping $\mathrm{Id} - t_f\nabla u_1(\cdot):\mathbb{R}^d \rightarrow \mathbb{R}^d$ is bijective and $\textbf{I}_d - t_f\nabla^2 u_1(x)\succ \textbf{O}_d$, then $\mathrm{Id} + t_f \nabla u_0(\cdot)$ is the optimal transport from $\mu$ to $\nu$, and $\mathrm{Id} - t_f \nabla u_1(\cdot)$ is the optimal transport map from $\nu$ to $\mu$.
\end{theorem}

\begin{proof}
We split the proof into several steps:

\textbf{Step 1.} We first prove the fact that $\nabla u_0(x - t_f\nabla u_1(x, t)) = \nabla u_1(x)$ for arbitrary $x\in\mathbb{R}^d$. This can be shown by taking gradient with respect to $x$ on both sides of \eqref{implicit_equ_at_T}:
\begin{equation*}
  (\textbf{I}_d - t_f\nabla^2 u_1(x)) \nabla u_0(x - t_f \nabla u_1(x)) = \nabla u_1(x) - t_f \nabla^2 u_1(x) \nabla u_1(x).
\end{equation*}
Re-arrange this equation yields
\begin{equation*}
  (\textbf{I}_d -t_f\nabla^2 u_1(x)) (\nabla u_0(x - t_f \nabla u_1(x)) - \nabla u_1(x)) = 0.
\end{equation*}
As $\textbf{I}_d-t_f\nabla^2 u_1(x)\succ \textbf{O}_d,$ we deduce that $\nabla u_0(x - t_f \nabla u_1(x)) = \nabla u_1(x)$ for arbitrary $x\in\mathbb{R}^d$.

\textbf{Step 2.} For the sake of brevity, we denote  $T_0(\cdot):=\mathrm{Id} + t_f \nabla u_0(\cdot)$, and $T_1(\cdot):=\mathrm{Id} - t_f\nabla u_1(\cdot)$. We prove that $T_0\circ T_1 = \mathrm{Id}$. This can be derived by straightforward calculation:
\begin{align*}
  T_0(T_1(x)) = T_1(x) + t_f\nabla u_0(T_1(x)) & =  x - t_f \nabla u_1(x) + t_f\nabla u_0(x - t_f \nabla u_1(x))  \\
  & = x + t_f(\nabla u_0(x - t_f \nabla u_1(x)) - \nabla u_1(x)) = x,
\end{align*}
for any $x \in \mathbb{R}^d$.

\textbf{Step 3.} Before we prove the assertion, we show that $\int_{\mathbb{R}^d} \|T_0(x)\|^2 \diff \mu < +\infty, \int_{\mathbb{R}^d} \|T_1(x)\|^2 \diff \nu < +\infty.$ The latter inequality can be shown using
\begin{align*}
  \int_{\mathbb{R}^d} \|T_1(x)\|^2 \diff \nu & \leq  2\left(\int_{\mathbb{R}^d} \|T_1(x) - x\|^2 \diff \nu + \int_{\mathbb{R}^d} \|x\|^2 \diff \nu\right) \\
  & =  2(t_f^2\|\nabla u_1\|_{L^2(\nu)}^2 + \int_{\mathbb{R}^d} \|x\|^2 \diff \nu) < +\infty.
\end{align*}
For the former inequality, we have
\begin{equation*}
  \int_{\mathbb{R}^d} \|T_0(x)\|^2 \diff \mu \leq % 2\left(\int_{\mathbb{R}^d} \|T_0(x) - x\|^2 \diff \mu + \int_{\mathbb{R}^d} \|x\|^2 \diff \mu\right) = 
  2 t_f^2\int_{\mathbb{R}^d} \|\nabla u_0(x)\|^2 \diff \mu + 2 \int_{\mathbb{R}^d} \|x\|^2 \diff \mu.
\end{equation*}
Using the fact that $T_{1\sharp}\nu = \mu,$ The first term above equals $\int_{\mathbb{R}^d} \|\nabla u_0(T_1(x))\|^2\diff \nu = \int_{\mathbb{R}^d} \|\nabla u_1(x)\|^2\diff \nu = \|\nabla u_1\|_{L^2(\nu)}<+\infty$, where we use the fact that $\nabla u_0(T_1(x)) = \nabla u_1(x)$ established in step 1. This accomplishes the proof for $\int_{\mathbb{R}^d} \|T_0(x)\|^2 \diff \mu < + \infty$.

\textbf{Step 4.} We now prove the conclusion. Firstly, recall that $D T_1(x) = \textbf{I}_d - t_f\nabla^2 u_1(x) \succ \textbf{O}_d,$ this leads to the fact that $T_1(\cdot)$ is strictly monotone. We now apply Lemma \ref{lemma:mono_map} to show that $T_0$ is also bijective, and strictly monotone. 

$T_0$ is bijective suggests that $T_1(\cdot)$ is the inverse mapping of $T_0(\cdot)$, this leads to $T_{1\sharp}\nu = \mu$. As $T_1(\cdot)=\nabla\varphi(\cdot)$ with $\varphi(\cdot) = \frac{\|\cdot\|^2}{2} - t_f u_1(\cdot)$ being convex, combinging the fact established in step 3, Lemma \ref{lemma:sufficient_cond_for_OT} proves that $T_1$ is the optimal transport map from $\nu$ to $\mu$. 

Furthermore, $T_0$ is strictly monotone yields that $DT_0(x) = \textbf{I}_d + t_f\nabla^2 u_0(x)\succ \textbf{O}_d$, this indicates that $T_0(\cdot)=\nabla\left(\frac{\|\cdot\|^2}{2} + t_f u_0(\cdot)\right)$ is the gradient of a convex function. Combining with the fact that $T_{0\sharp}\mu = \nu$ and finite $L^2(\mu)$ cost for $T_0$, we deduce that $T_0$ is the optimal transport map form $\mu$ to $\nu$.
\end{proof}

\end{document}
